%% kapa: exam=1 round=20 year=2009 period=1 form=A subject=economics subject_ko=경제원론 seq=14 number=54
%% source: Q-Net 감정평가사 기출문제 (cst00309) artlSeq=5024178 file=extracted/1008/5024178/2009년 제20회 감정평가사 1차시험 기출문제/1교시.hwp
%% conversion: kordoc 4.24.51 -> tools/gen1.py (auto); review=proofread-sonnet; flags=none
\begin{question}{54}
소비자 A는 사과와 귤만을 소비한다. 주어진 예산을 다 쓸 경우 사과 10개와 귤 20개를 살 수도 있고, 사과 14개와 귤 16개를 살 수도 있다. 만약 주어진 예산을 모두 사과만을 사는데 쓴다면 살 수 있는 사과의 갯수는? (단, 사과와 귤의 가격은 일정하다.)
\begin{choices}
  \item 20개
  \item 24개
  \item 30개
  \item 36개
  \item 40개
\end{choices}
\end{question}
\answer{3}
